% =====================================================================
%  JEPA-IR — Slides de présentation (Beamer)
%  Compile avec :  pdflatex slides.tex  (2x pour la table des matières)
% =====================================================================
\documentclass[aspectratio=169,11pt]{beamer}

% --- Thème : NATIF Beamer (compile partout, y compris Overleaf sans paquet sup.)
%     Si tu as 'metropolis' installé, tu peux remplacer les 2 lignes ci-dessous
%     par : \usetheme{metropolis}
\usetheme{default}
\usecolortheme{whale}
\useinnertheme{circles}

\usepackage[utf8]{inputenc}
\usepackage[T1]{fontenc}
\usepackage[french]{babel}
\usepackage{booktabs}
\usepackage{tikz}
\usetikzlibrary{arrows.meta,positioning,fit,backgrounds}
\usepackage{xcolor}

% --- palette ---
\definecolor{jblue}{HTML}{1f6feb}
\definecolor{jgreen}{HTML}{2ca02c}
\definecolor{jred}{HTML}{d62728}
\definecolor{jorange}{HTML}{d29922}
\definecolor{jgrey}{HTML}{8b949e}

\setbeamertemplate{navigation symbols}{}
\setbeamercolor{frametitle}{bg=jblue,fg=white}
\setbeamercolor{title}{fg=jblue}
\setbeamercolor{structure}{fg=jblue}
\setbeamertemplate{blocks}[rounded][shadow=true]

\newcommand{\hl}[1]{\textcolor{jblue}{\textbf{#1}}}
\newcommand{\good}[1]{\textcolor{jgreen}{\textbf{#1}}}
\newcommand{\bad}[1]{\textcolor{jred}{\textbf{#1}}}

\title{JEPA-IR}
\subtitle{Encodeur auto-supervisé de programmes \\ \& prédiction latente de l'optimisation compilateur}
\author{Hack The World(s) — équipe \texttt{jepacode}}
\date{\today}

\begin{document}

\maketitle

\begin{frame}{Plan}
  \tableofcontents
\end{frame}

% =====================================================================
\section{Le problème}
% =====================================================================
\begin{frame}{Optimiser du code : un problème coûteux}
  \begin{block}{Constat}
    Pour savoir si une optimisation accélère un programme, il faut aujourd'hui
    le \bad{compiler ET l'exécuter} — répété sur des milliers de variantes
    (\emph{phase-ordering}). Prohibitif à grande échelle.
  \end{block}
  \pause
  \begin{alertblock}{Idée}
    Et si on \hl{raisonnait sur une représentation numérique du code},
    \hl{sans l'exécuter} ?
  \end{alertblock}
  \pause
  \begin{itemize}
    \item Transformer un programme en \textbf{vecteur} qui capture sa sémantique.
    \item Prédire l'effet de l'optimisation \textbf{dans l'espace latent}.
  \end{itemize}
\end{frame}

% =====================================================================
\section{Le pipeline}
% =====================================================================
\begin{frame}{Vue d'ensemble du pipeline}
  \centering
  \begin{tikzpicture}[
      node distance=8mm and 10mm,
      box/.style={draw, rounded corners, align=center, minimum height=11mm,
                  font=\small, fill=gray!8},
      arr/.style={-{Stealth[length=3mm]}, thick, jblue}]
    \node[box] (code) {code\\C/C++};
    \node[box, right=of code] (ir) {LLVM IR};
    \node[box, right=of ir, fill=jblue!12] (graph) {graphe\\3-relations};
    \node[box, right=of graph, fill=jgreen!15] (enc) {encodeur\\GNN (JEPA)};
    \node[box, right=of enc] (emb) {embedding};
    \draw[arr] (code)--(ir) node[midway,above,font=\tiny]{clang};
    \draw[arr] (ir)--(graph);
    \draw[arr] (graph)--(enc);
    \draw[arr] (enc)--(emb);
    \node[box, below=14mm of enc, fill=jorange!15] (pred) {predictor\\latent};
    \draw[arr] (emb.south) |- (pred.east);
    \node[font=\scriptsize, jgrey, below=1mm of pred] {emb(O$_k$) $\to$ emb(O$_{k+1}$)};
  \end{tikzpicture}
  \vspace{4mm}
  \begin{itemize}\footnotesize
    \item \hl{Étape clang} : locale (Dalia n'a pas de clang).
    \item \hl{Encodeur} : entraîné \emph{de zéro}, auto-supervisé, sur GPU B200.
  \end{itemize}
\end{frame}

% =====================================================================
\section{La représentation : graphe 3-relations}
% =====================================================================
\begin{frame}{Un graphe d'IR \emph{sans perte} : 3 relations}
  Chaque \textbf{nœud} = une instruction LLVM. Trois types d'\textbf{arêtes},
  conservés simultanément :
  \vspace{2mm}
  \begin{itemize}
    \item \textcolor{jred}{\textbf{contrôle}} — flot de contrôle (branches du CFG)
    \item \textcolor{jblue}{\textbf{données}} — flot de données (def$\to$use SSA)
    \item \textcolor{jgreen}{\textbf{mémoire}} — ordre des effets de bord (load/store)
  \end{itemize}
  \vspace{3mm}
  \begin{block}{Ce qui le distingue}
    Ni un AST, ni un simple CFG, ni ProGraML (control+data+call) : on garde
    \hl{l'ordre des effets mémoire} — aucune des trois relations n'est jetée.
  \end{block}
\end{frame}

% =====================================================================
\section{L'encodeur JEPA}
% =====================================================================
\begin{frame}{JEPA : apprendre une représentation, sans labels}
  \begin{columns}[T]
    \begin{column}{0.54\textwidth}
      \textbf{Principe (sans décodeur) :}
      \begin{enumerate}\small
        \item masquer une partie du graphe (nœuds + arêtes)
        \item encoder le graphe \textbf{masqué} $\to z_A$
        \item encoder le graphe \textbf{complet} $\to z_B$
        \item rapprocher $z_A$ et $z_B$ \hl{dans l'espace latent}
      \end{enumerate}
      \vspace{2mm}
      \footnotesize La cible est un \textbf{vecteur}, jamais le graphe.
      On ne reconstruit rien.
    \end{column}
    \begin{column}{0.44\textwidth}
      \begin{block}{Anti-collapse (obligatoire)}
        \footnotesize
        Sans contrainte, l'encodeur peut tricher : sortir un vecteur constant.
        \\ \hl{VICReg} (variance + covariance) l'interdit.
      \end{block}
      \begin{exampleblock}{Loss}
        \scriptsize
        $\mathcal{L} = \underbrace{\mathrm{MSE}(z_A,z_B)}_{\text{invariance}}
          + \underbrace{\mathrm{var}}_{\text{anti-collapse}}
          + \underbrace{\mathrm{cov}}_{\text{décorrélation}}$
      \end{exampleblock}
    \end{column}
  \end{columns}
\end{frame}

\begin{frame}{L'encodeur ne s'effondre pas (preuve par ablation)}
  \begin{columns}[T]
    \begin{column}{0.5\textwidth}
      \textbf{Avec VICReg :}
      \begin{itemize}\small
        \item std des embeddings $\approx \good{1.3}$ (stable)
        \item rang effectif $\approx 30/128$
        \item nuage PCA étalé, pas de point/ligne
      \end{itemize}
    \end{column}
    \begin{column}{0.5\textwidth}
      \textbf{Ablation — sans VICReg :}
      \begin{itemize}\small
        \item std $\to \bad{0.0003}$ (effondrement total)
        \item l'encodeur sort un vecteur constant
      \end{itemize}
    \end{column}
  \end{columns}
  \vspace{4mm}
  \begin{center}
    \begin{tikzpicture}
      \node {std embeddings :};
      \node[right=2mm, jgreen, font=\large\bfseries] {1.3};
      \node[right=24mm, font=\large] {vs};
      \node[right=38mm, jred, font=\large\bfseries] {0.0003};
    \end{tikzpicture}
  \end{center}
  \begin{center}\small Preuve par contraste : la régularisation est ce qui empêche le collapse.\end{center}
\end{frame}

% =====================================================================
\section{Le predictor d'optimisation}
% =====================================================================
\begin{frame}{Predictor : prédire l'optimisation dans le latent}
  Un petit MLP apprend à \hl{avancer d'un cran d'optimisation} :
  \[ \text{predictor}\big(\text{emb}(O_k)\big) \approx \text{emb}(O_{k+1}) \]
  \begin{itemize}
    \item Travaille \textbf{uniquement sur les vecteurs} (pas de génération de code).
    \item Encodeur \textbf{gelé} ; seul le MLP est entraîné (loss = MSE latente).
    \item Évalué sur un \hl{test set jamais vu} (pools disjoints, anti-fuite).
  \end{itemize}
  \vspace{2mm}
  \begin{block}{Méthodologie anti-fuite}
    \footnotesize Corpus découpé en 3 pools disjoints par hash
    (\texttt{encoder} / \texttt{predictor} / \texttt{heldout}), split
    train/val/test 70/15/15. L'encodeur n'a jamais vu les programmes du predictor.
  \end{block}
\end{frame}

\begin{frame}{Résultat : le predictor capture le saut O0$\to$O1}
  \centering
  \small
  \begin{tabular}{lccccc}
    \toprule
    Transition & MSE id. & MSE pred. & gain & cos(in,cible) & cos(pred,cible) \\
    \midrule
    \textbf{O0$\to$O1} & 0.912 & 0.393 & \good{+56.9\%} & 0.795 & \good{0.915} \\
    O1$\to$O2 & 0.094 & 0.085 & +9.1\% & 0.982 & 0.983 \\
    O2$\to$O3 & 0.018 & 0.025 & \bad{$-$36.3\%} & 0.997 & 0.995 \\
    \bottomrule
  \end{tabular}
  \vspace{4mm}
  \begin{itemize}\small
    \item \good{O0$\to$O1} (le vrai saut) : le predictor reconstruit l'embedding
          optimisé, \hl{cos 0.79 $\to$ 0.92}.
    \item O1, O2, O3 \emph{quasi confondus} $\Rightarrow$ presque rien à prédire.
  \end{itemize}
\end{frame}

% =====================================================================
\section{Limites \& honnêteté scientifique}
% =====================================================================
\begin{frame}{Pourquoi O1 $\approx$ O2 $\approx$ O3 ?}
  \begin{block}{Mesure (800 programmes, IR brut O2 vs O3)}
    \centering \hl{IR identique à 100\,\%} — clang ne change \emph{rien} en -O3.
  \end{block}
  \begin{itemize}
    \item \textbf{La vraie limite = le corpus.} AnghaBench = fonctions
          \emph{isolées et courtes} ($\sim$22 nœuds).
    \item -O3 cible l'inlining inter-procédural et la vectorisation
          $\Rightarrow$ \hl{rien à optimiser} sur du code isolé.
    \item Ce n'est pas un défaut de la représentation : l'IR \emph{lui-même}
          est identique (0\,\% de perte d'info mesurée).
  \end{itemize}
  \vspace{1mm}
  \footnotesize\textit{Hiérarchie de saturation : O0$\to$O1 gros saut ($\sim$89\,\% des graphes changent),
  O1$\to$O2 petit ($\sim$16\,\%), O2$\to$O3 nul (0\,\%).}
\end{frame}

\begin{frame}{Honnêteté scientifique}
  \begin{itemize}
    \item L'IR est \textbf{déterministe} : le JEPA apprend une \emph{représentation}
          (comme BERT sur du texte), pas de l'aléatoire. La vraie prédiction =
          deviner emb($O_{k+1}$) sans compiler.
    \item Pas de \textbf{temps d'exécution réel} (corpus non exécutable)
          $\Rightarrow$ on prédit l'effet \emph{structurel}, pas le speedup.
    \item \hl{Prochaine étape} : corpus de programmes complets et exécutables
          (ExeBench) où -O3 fait une vraie différence.
  \end{itemize}
\end{frame}

% =====================================================================
\section{Applications}
% =====================================================================
\begin{frame}{Applications}
  \centering\small
  \begin{tabular}{ll}
    \toprule
    Usage & Apport \\
    \midrule
    Phase-ordering & scorer des séquences d'optim sans les exécuter \\
    Auto-tuning compilateur & choisir l'optim adaptée à \emph{ce} code \\
    Repr. réutilisable & embedding « universel » de code machine \\
    Similarité / clones & détecter code copié, malware, vulnérabilités \\
    Code search & chercher du code par le \emph{sens} \\
    \bottomrule
  \end{tabular}
\end{frame}

% =====================================================================
\begin{frame}
  \vfill
  \centering
  {\Huge\textcolor{jblue}{\textbf{JEPA-IR}}}\\[4mm]
  {\large Un encodeur de programmes qui « comprend » le code,\\
  et prédit l'optimisation sans l'exécuter.}
  \vfill
\end{frame}

\begin{frame}{Récapitulatif}
  \begin{itemize}
    \item \good{Encodeur JEPA} : représentation 3-relations, aucun collapse
          (prouvé par ablation).
    \item \good{Predictor latent} : capture le saut O0$\to$O1 (cos 0.79$\to$0.92).
    \item \good{Méthodologie} : pools disjoints, train/val/test, zéro fuite.
    \item \good{Limite mesurée} : clang sature à -O2 sur ce corpus
          $\Rightarrow$ étendre à des programmes complets.
  \end{itemize}
\end{frame}

\end{document}